\scanpage{149}
passent à la limite.

Soit \deleted{$\unsure{Z}$} $\unsure{D|Z}$ un ens. Une graduation de type $D$ de $M$ est une décomposition en somme directe
\[
M\simeq\bigoplus_{d\in\unsure{D|Z}}M_d\qquad(M_d\text{ sous-module de }M)
\]
\deleted{\unsure{On interprète}}\inserted{\unsure{s'écrit}} aussi comme décomposition de $V(M)$ en produit
\[
V(M)\simeq\prod_{d\in\unsure{D|Z}}V(M_d)
\]
(modules quotients). Il est souvent commode aussi de regarder \deleted{la cat.}\inserted{(\unsure{cat. oppos. de cat.})} des \inserted{$M_*$} gradués par $D$ comme des familles $(M_d)_{d\in D}$ de $k$-modules, donc $V(M)_{\inserted{*}}$ comme la \uline{famille} des $V(M_d)$. $(d\in Z)$

Si $\unsure{D|Z}$ est muni d'une structure de groupe abélien, alors la structure dite sur $M$ revient à faire opérer le groupe diagonalisable de $\widehat C_k$
\diagram{
\begin{tikzcd}[column sep=large]
D(Z)=G: \arrow[d,phantom,"\wr" description] & A\arrow[r,phantom,"\unsure{-|\longmapsto}" description] & \operatorname{Hom}_{\mathrm{Ab}}(\unsure{D|Z},A^*)\arrow[d,phantom,"\wr" description]\\
\operatorname{Spec}k(Z) & & \operatorname{Hom}_{k\text{-Alg}}(\underbrace{k(Z)}_{\text{alg du groupe}},A)
\end{tikzcd}
}
\inserted{lin. sur}\deleted{sur} $M$ [\unsure{par quoi on entend} : sur $W(M):A\mathbin{\text{---}}M\otimes A$, ou au choix par
\margin{148}
